\section{Performance Evaluation}
\label{sec:evaluation}
In this section, we conduct extensive experiments to answer the following research question: \textbf{Which SSH Connection Protocol can meet the strict latency and completeness demands of closed-loop terminal agents?}

\subsection{Experiment Setup}
\label{subsec:experiment-setup}

\sout{Figure~\ref{fig:setup} shows the controlled LAN setup used to evaluate SSHv2, SSH3, and Mosh under three emulated network profiles, as summarized in Table~\ref{tab:network-scenarios}. These profiles are configured using Linux tc/netem to create a controlled range of network quality. This allows us to observe how each protocol behaves as network impairment increases, and whether its performance degrades gradually or drops sharply under more severe conditions.}
\textcolor{blue}{Figure~\ref{fig:setup} shows the experimental setup used to evaluate SSHv2, SSH3, and Mosh across both the controlled LAN and VPN-based environments. In the controlled LAN experiments, we use three emulated network profiles, as summarized in Table~\ref{tab:network-scenarios}. These profiles are configured using Linux \texttt{tc/netem} to provide a controlled range of delay, jitter, and packet loss. This setup allows us to examine how the performance of each protocol changes as network impairment increases, and whether its performance degrades gradually or drops sharply under more severe conditions.}

\sout{Figure~\ref{fig:setup} shows the VPN-based Internet-path setup. Unlike the controlled LAN experiments, the VPN scenario uses measured path characteristics collected during the experiments rather than preconfigured impairments. This setup complements the controlled profiles by evaluating the protocols over a real network path with naturally varying latency, jitter, and loss.}
\textcolor{blue}{In the VPN experiment, the benchmark client and remote server remain the same, but traffic is routed through a VPN over the Internet instead of the controlled LAN scenario. Unlike the LAN scenarios, we do not apply fixed delay, jitter, or packet loss settings with \texttt{tc/netem}; instead, we record the network conditions observed during the experiment. This VPN scenario complements the controlled LAN scenarios by showing how each protocol behaves under naturally varying Internet latency, jitter, and loss.}

\sout{The experiments use one VPN-based Internet path and three emulated LAN profiles, as summarized in Table~\ref{tab:network-scenarios}.}

\sout{The three emulated profiles are selected to create a controlled range of path quality, allowing us to observe how each protocol behaves as network conditions become increasingly impaired.}

\sout{This design highlights whether protocol behavior degrades smoothly or changes sharply under stronger impairment.}

\sout{The profiles are configured with Linux tc/netem to control the emulated path conditions.}

\sout{The VPN scenario uses measured path characteristics collected during the experiments rather than configured impairments.}




\begin{table}[t]
\centering
\caption{Network configuration parameters in the evaluated network environments}
\label{tab:network-scenarios}

\resizebox{\columnwidth}{!}{%
\begin{tabular}{llccc}
\hline
\multicolumn{1}{c}{\multirow{2}{*}{Network}} &
\multicolumn{1}{c}{\multirow{2}{*}{Dynamicity level}} &
OWD & Jitter & Loss \\
& & (ms) & (\%) & (\%) \\
\hline

\multirow{3}{*}{Simulated}
    & Low    & 10.0  & 0.0  & 0.0 \\
    & Medium & 50.0  & 4.0  & 1.5 \\
    & High   & 100.0 & 16.0 & 3.0 \\

\hline
VPN & Based on real-network conditions & 78.2 & 11.3 & 0.0 \\
\hline
\end{tabular}%

}
\end{table}

The setup consists of a benchmark client, a network path, and a remote server. Two Raspberry Pi 4 devices are used as the client and the server, respectively. Because Linux \texttt{netem} \sout{shapes egress traffic} \textcolor{blue}{only applies network attenuation to outgoing traffic}, the same \sout{profile} \textcolor{blue}{scenario} is applied on both hosts to emulate bidirectional degradation. A configured one-way delay $D$ therefore yields approximately $2D$ ms RTT, and a per-egress loss rate $p$ yields approximately $1-(1-p)^2$ round-trip loss exposure. We evaluate concrete protocol implementations rather than abstract protocol models. SSHv2 is implemented with OpenSSH 9.2 and OpenSSL 3.0.17. Mosh is implemented with Mosh 1.4.0 and uses the standard SSH bootstrap procedure before switching to its UDP-based State Synchronization Protocol channel. SSH3 is implemented using the francoismichel/ssh3 implementation and connects to the SSH3 terminal server endpoint \cite{michel2023}. All protocols use the same remote account, SSH key material, server host, terminal environment, and workload content.

For each protocol, workload, and scenario, the benchmark runs 50 independent trials. Each trial establishes a fresh session and then executes 100 commands. Session setup latency is measured once per trial, giving 10 setup measurements for each configuration. Workload latency is measured once per command or probe iteration, giving 1000 workload samples for each configuration. Within each scenario, protocol, workload, and command or probe combinations are shuffled \sout {using a fixed seed} \textcolor{blue}{using a fixed initialization value} to reduce order effects from cache state, device temperature, run order, and background variation.

\begin{figure}[t]
\centering
\includegraphics[width=0.9\columnwidth]{figs/setup.pdf}
\caption{Experimental setup for evaluating SSHv2, SSH3, and Mosh across controlled LAN and VPN-based environments.}
\label{fig:setup}
\end{figure}

\subsection{Experiment Results}
\label{subsec:experiment-results}

\begin{figure}[h]
\centering

\makebox[\columnwidth][c]{%
    \hspace*{0.05\columnwidth}%
    
\includegraphics[width=0.59\columnwidth]
    {figs/session_setup_legend_cropped.pdf}%
    \hspace*{-0.05\columnwidth}%
}\\[1.5ex]

\includegraphics[width=\columnwidth]
{figs/session_setup_bar.pdf}

\caption{Mean session setup latency across four network scenarios.}
\label{fig:session-setup}

\end{figure}


This section evaluates SSHv2, Mosh, and SSH3 \sout{across session setup}\textcolor{blue}{in terms of session setup latency}, \sout{workloads latency}\textcolor{blue}{workload latency}, and output completeness. Figure~\ref{fig:session-setup} reports \sout{setup latency,}\textcolor{blue}{session setup latency.} \sout{Figures}\textcolor{blue}{Figure}~\ref{fig:w1-w2-performance} \sout{and \ref{fig:w3-performance} report steady-state workload performance, and}\textcolor{blue}{reports the results for W1 and W2, which represent command workloads with small and large outputs, respectively, together with their output completeness. Figure~\ref{fig:w3-performance} reports the results for W3, which evaluates interactive editing latency under single pane and five pane terminal layouts.} \sout{A table gives the exact output-completeness values for W1 and W2.} We \sout{present the results by measurement category and explain}\textcolor{blue}{use these workload results to explain} the observed differences using the protocol behavior discussed earlier.




\begin{figure*}[t]
    \centering
    
\includegraphics[width=0.70\textwidth]{figs/latency_workload_legend.pdf}\\[-0.4ex]
    \subfloat[W1: Small-output latency and output completeness]{%
        \includegraphics[width=0.48\textwidth]{figs/latency_completeness_w1_by_scenario.pdf}%
        \label{fig:latency-w1}
    }\hfill
    \subfloat[W2: Large-output latency and output completeness]{%
        \includegraphics[width=0.48\textwidth]{figs/latency_completeness_w2_by_scenario.pdf}%
        \label{fig:latency-w2}
    }
    \caption{Mean latency and 95\% confidence intervals for W1 and W2, excluding warmup samples. Bars show workload latency, while markers show output completeness as the percentage of expected bytes received.}
    \label{fig:w1-w2-performance}
\end{figure*}

\begin{figure*}[t]
    \centering
    
\includegraphics[width=0.70\textwidth]{figs/latency_workload_legend.pdf}\\[-0.4ex]
    \subfloat[W3-1p: Single-pane]{%
        \includegraphics[width=0.48\textwidth]{figs/latency_w3_1p_by_scenario.pdf}%
        \label{fig:latency-w3-1p}
    }\hfill
    \subfloat[W3-5p: Five-pane]{%
        \includegraphics[width=0.48\textwidth]{figs/latency_w3_5p_by_scenario.pdf}%
        \label{fig:latency-w3-5p}
    }
    \caption{Mean keystroke-to-echo latency and 95\% confidence intervals for the W3 interactive-editing workload under single-pane and five-pane terminal layouts.}
    \label{fig:w3-performance}
\end{figure*}
\captionsetup{font=normalsize}
\captionsetup[subfloat]{font=normalsize}

Session setup latency is measured from client process spawn to the first interactive shell prompt. Figure~\ref{fig:session-setup} shows that SSH3 has the lowest setup latency in every network scenario. \sout{On the VPN path, SSH3 reaches the first usable prompt in 876.0\,ms, compared with 2873.3\,ms for SSHv2 and 3280.0\,ms for Mosh. The same ordering holds under emulated impairment: SSH3 requires 982.3\,ms in Low, 1294.2\,ms in Medium, and 1685.7\,ms in High, while SSHv2 requires 1674.3\,ms, 2577.1\,ms, and 3960.5\,ms, and Mosh requires 2270.5\,ms, 3503.6\,ms, and 5161.6\,ms} \textcolor{blue}{In the VPN scenario, SSH3 took 876.0\,ms to reach the first usable shell prompt, compared with 2873.3\,ms for SSHv2 and 3280.0\,ms for Mosh. The same ordering holds in the controlled LAN scenarios: SSH3 took 982.3\,ms, 1294.2\,ms, and 1685.7\,ms to reach the first usable shell prompt in the Low, Medium, and High scenarios, respectively. For comparison, SSHv2 took 1674.3\,ms, 2577.1\,ms, and 3960.5\,ms, while Mosh took 2270.5\,ms, 3503.6\,ms, and 5161.6\,ms under the same scenarios}. This result follows from the connection-establishment path of each protocol. SSH3 runs over HTTP/3 and QUIC \cite{michel2023}, whose handshake combines transport and cryptographic negotiation \cite{Iyengar2021}. SSHv2 first establishes TCP and then performs SSH key exchange and authentication \cite{Ylonen2006RFC4253}. \sout{Mosh pays the highest startup cost because it uses SSHv2 to authenticate and bootstrap the session before switching to its UDP-based State Synchronization Protocol channel \cite{Winstein2012}. For an AI-like terminal client, this cold-start cost matters when the client opens fresh shells, reconnects after failures, or runs short-lived remote-control loops}. \textcolor{blue}{Mosh had the longest session setup time in these measurements because it uses SSHv2 for authentication and session establishment before switching to a UDP-based State Synchronization Protocol (SSP) channel \cite{Winstein2012}. In the context of interactive terminal tasks like AI, this session startup latency can impact measurement results and execution efficiency, especially when the task frequently opens new shells, restores connections, or runs multiple short remote control loops}.

W1 measures simulated interactions on commands with short response times and small output. As shown in Figure~\ref{fig:latency-w1}, SSH3 \sout{achieved} \textcolor{blue}{achieved} the lowest mean W1 latency across four network scenarios, although the margins were small on slightly degraded paths. In a real-world internet environment, SSH3 completed the command loop in 177.9 ms, compared to 188.8 ms for SSHv2 and 190.9 ms for Mosh. In the Low volatility network scenario, SSH3 and SSHv2 were nearly equal at 75.5 ms and 76.1 ms, respectively, while Mosh achieved 82.2 ms. The difference became apparent when network volatility was at High, where SSH3 achieved 307.5 ms, outperforming Mosh at 317.3 ms and SSHv2 at 367.0 ms. This result shows that the QUIC-based terminal channel of SSH3 handles short response loops efficiently. QUIC integrates encryption, encapsulation, acknowledgment, and loss recovery at the transport layer. In SSH3, this gives the tested terminal channel a shorter W1 completion path compared to SSHv2 over TCP and Mosh over SSP. SSHv2 preserves the completeness of the byte stream, but the loss and recovery of the stream are delayed, increasing completion latency in faulty paths. Mosh is optimized for more intuitive responsiveness rather than full byte command loop completion, so it offers less advantage to W1 \sout{as shown in the table}\textcolor{blue}{as reflected by the W1 latency and output-completeness results in Figure~\ref{fig:latency-w1}}.

W2 \sout{stresses byte-complete delivery}\textcolor{blue}{emphasizes the full transmission of every byte} of large terminal output. Mosh \sout{appears fastest, but this latency is not comparable: Table shows that}\textcolor{blue}{has the fastest latency; however, the W2 output-latency and output-completeness results in Figure~\ref{fig:latency-w2} indicate that} it \sout{receives}\textcolor{blue}{only receives} 1.43\%, 4.19\%, 1.75\%, and 1.24\% of the expected W2 bytes \sout{across VPN, Low, Medium, and High}\textcolor{blue}{in the VPN, Low, Medium, and High network scenarios, respectively} because \sout{SSP synchronizes screen state}\textcolor{blue}{the protocol's SSP synchronizes terminal screen state} rather than \sout{every}\textcolor{blue}{each} intermediate byte \cite{Winstein2012}. Among \sout{complete-output transports}\textcolor{blue}{the methods for full output transmission}, SSHv2 is consistently faster than SSH3: for example, SSHv2 completes W2 in 4213.5 ms \sout{versus}\textcolor{blue}{compared to} 47732.6 ms for SSH3 over VPN, and 36363.8 ms \sout{versus}\textcolor{blue}{compared to} 74664.5 ms \sout{under High impairment}\textcolor{blue}{in a High-level dynamic network scenario}. The main reason is that W2 uses \sout{one terminal session}\textcolor{blue}{terminal sessions} and sends the command output sequentially through \sout{one}\textcolor{blue}{a} shell stream. QUIC avoids head-of-line blocking across independent streams, but bytes within a single stream are still \sout{delivered}\textcolor{blue}{transmitted} in order. Therefore, \sout{when}\textcolor{blue}{because} W2 uses only one stream, SSH3 does not benefit from QUIC's \sout{multi-stream recovery behavior}\textcolor{blue}{multi-stream recovery behavior} \cite{Iyengar2021}. In this \sout{single-stream regime}\textcolor{blue}{single-stream mode}, both SSHv2 and SSH3 \sout{are exposed to ordered-delivery stalls}\textcolor{blue}{experience bottlenecks due to sequential transmission}. SSHv2 benefits from the \sout{mature TCP path implemented and optimized in the kernel}\textcolor{blue}{TCP path already implemented and optimized in the operating system kernel}. \sout{By contrast}\textcolor{blue}{Conversely}, QUIC stacks are \sout{commonly implemented in the user space over UDP, which adds packet processing}\textcolor{blue}{often implemented in user space via UDP, which adds to packet processing}. This result is consistent with recent implementation results showing that QUIC throughput depends on application \sout{use}\textcolor{blue}{usage}, stream concurrency, implementation behavior, and packet processing overhead \cite{Konig2025}.

W3 evaluates keystroke-to-echo latency to capture client-visible responsiveness during interactive terminal use. Overall, Mosh consistently outperforms the two byte-stream protocols (SSH3 and SSHv2). In the 1p \texttt{tmux} layout in Figure~\ref{fig:latency-w3-1p}, Mosh demonstrates a substantial advantage on impaired paths, maintaining a latency of just 25.3 ms under High impairment. In contrast, both SSHv2 and SSH3 experience severe degradation, with both exceeding 220 ms under the same conditions. This behavioral gap in the 1p scenario stems from fundamental architectural differences. Mosh achieves Low latency through predictive local echo and screen-state synchronization, hiding network delay from client-visible rendering \cite{Winstein2012}. Conversely, both SSH3 and SSHv2 operate strictly as end-to-end byte streams. Regardless of whether they utilize modern QUIC or legacy TCP, both protocols must wait for the remote server to process the keystroke and send the output back, making their responsiveness entirely dependent on network round-trip time and terminal rendering overhead. Moving to the more complex 5p layout in Figure~\ref{fig:latency-w3-5p}, the performance gap between Mosh and the byte-stream protocols widens further. Mosh remains highly responsive, recording only 31.3 ms under High impairment. Meanwhile, the latency for both SSH3 and SSHv2 scales up significantly. This magnified delay occurs because increasing the number of panes increases the volume of screen state updates that the server must push to the client. Since SSH3 and SSHv2 are obligated to transmit every byte of the terminal redrawing sequence, their interactive latency inherently inflates. In contrast, Mosh efficiently coalesces these complex layout updates and predicts visible changes locally, preventing the network from becoming a rendering bottleneck. Consequently, W3 represents a fundamentally different operational regime compared to W1 and W2. While SSH3 and SSHv2 exhibit restrictive, network-bound characteristics during interactive editing, Mosh stands out as the definitively superior choice for masking network impairments and preserving human-driven responsiveness.


Returning to our central research question, our evaluation reveals that no single transport universally satisfies all phases of agent interaction. Instead, the optimal protocol is workload-dependent. For agile, short-lived interactions requiring rapid session establishment and Low command-response latency, SSH3’s QUIC-based transport provides the most efficient feedback channel. However, for programmatic observation phases that rely on parsing large terminal outputs, legacy SSHv2 remains the superior choice, delivering the fastest byte-complete transmission despite its TCP roots. Crucially, our results demonstrate that transport mechanisms optimized for human visual responsiveness, such as Mosh, inherently conflict with the programmatic needs of automated clients, sacrificing critical byte-completeness in heavy-output scenarios.
